\heading{实验物理中的统计方法-模拟题3-期末}
\noteinfo{题目和答案由AI生成。知识点范围按现有往年题整理，叙述风格尽量向往年题靠拢；本套难度较前两套期末模拟题再提高一些。}

\begin{question}
设物理量
\[
W=\ln\frac{X}{Y},
\]
其中 $X,Y$ 的标准差分别为 $\sigma_X,\sigma_Y$，协方差为 $c=\operatorname{Cov}(X,Y)$。写出小误差近似下 $W$ 的不确定度公式。
\begin{solution}
由误差传播公式，
\[
\sigma_W^2\approx
\left(\frac{\partial W}{\partial X}\right)^2\sigma_X^2+
\left(\frac{\partial W}{\partial Y}\right)^2\sigma_Y^2+
2\frac{\partial W}{\partial X}\frac{\partial W}{\partial Y}c.
\]
而
\[
\frac{\partial W}{\partial X}=\frac1X,
\qquad
\frac{\partial W}{\partial Y}=-\frac1Y.
\]
故
\ansmath{\sigma_W^2\approx \frac{\sigma_X^2}{X^2}+\frac{\sigma_Y^2}{Y^2}-\frac{2c}{XY}}
\end{solution}
\end{question}

\begin{question}
设 $X_1,\cdots,X_n$ 独立同分布，且概率密度为
\[
f(x;\theta)=\frac{2x}{\theta^2},\qquad 0<x<\theta,
\]
其中 $\theta>0$。
\begin{enumerate}[label=(\alph*),leftmargin=2.8em,itemsep=0.2em]
\item 求 $\theta$ 的矩估计量；
\item 求 $\theta$ 的极大似然估计量；
\item 由样本最大值构造一个无偏估计量。
\end{enumerate}
\begin{solution}
先求总体均值：
\[
E[X]=\int_0^{\theta} x\cdot \frac{2x}{\theta^2}\,dx
=\frac{2\theta}{3}.
\]
故矩估计量满足
\[
\bar X=\frac{2\theta}{3},
\]
所以
\ansmath{\hat\theta_{\rm MOM}=\frac{3}{2}\bar X}
似然函数为
\[
L(\theta)=\prod_{i=1}^n \frac{2x_i}{\theta^2}\,\mathbf 1(0<x_i<\theta)
\propto \theta^{-2n}\mathbf 1(\theta\ge X_{(n)}),
\]
其中 $X_{(n)}=\max(X_1,\cdots,X_n)$。显然当 $\theta$ 取到允许范围内最小值时似然最大，因此
\ansmath{\hat\theta_{\rm MLE}=X_{(n)}}
又因为
\[
P(X_{(n)}\le x)=\left(\frac{x^2}{\theta^2}\right)^n,
\qquad 0<x<\theta,
\]
故其密度为
\[
f_{X_{(n)}}(x)=\frac{2n}{\theta^{2n}}x^{2n-1},
\qquad 0<x<\theta.
\]
于是
\[
E[X_{(n)}]=\frac{2n}{2n+1}\theta.
\]
因此无偏估计量可取
\ansmath{\tilde\theta=\frac{2n+1}{2n}X_{(n)}}
\end{solution}
\end{question}

\begin{question}
某正态总体方差已知为 $\sigma^2=4$。现从该总体中抽取 $n=16$ 个样本，样本均值记为 $\bar X$。考虑检验
\[
H_0:\mu=0,
\qquad
H_1:\mu>0.
\]
若采用拒绝域 $\bar X>c$，要求显著性水平为 $\alpha=0.05$。
\begin{enumerate}[label=(\alph*),leftmargin=2.8em,itemsep=0.2em]
\item 求临界值 $c$；
\item 当真实均值 $\mu=1$ 时，求检验功效。
\end{enumerate}
可取 $z_{0.95}=1.645$。
\begin{solution}
在 $H_0$ 下，
\[
\bar X\sim N\left(0,\frac{4}{16}\right)=N\left(0,\frac14\right).
\]
要求
\[
P_{H_0}(\bar X>c)=0.05,
\]
故
\[
c=z_{0.95}\times \frac{2}{4}=1.645\times 0.5=0.8225.
\]
所以
\ansmath{c=0.8225}
当 $\mu=1$ 时，
\[
\bar X\sim N\left(1,\frac14\right).
\]
故功效为
\[
P(\bar X>0.8225\mid \mu=1)
=P\left(Z>\frac{0.8225-1}{0.5}\right)
=P(Z>-0.355).
\]
数值上
\[
P(Z>-0.355)\approx 0.639.
\]
因此
\ansmath{\text{power at }\mu=1\approx 0.639}
\end{solution}
\end{question}

\begin{question}
某物理量有两次测量结果：
\[
A_1=1.20\pm 0.10,
\qquad
A_2=1.00\pm 0.20,
\]
二者的相关系数为 $\rho=0.2$。设要把这两个结果合并成一个最优线性无偏估计 $\hat A=w_1A_1+w_2A_2$（其中 $w_1+w_2=1$），求 $\hat A$ 及其不确定度。
\begin{solution}
协方差矩阵为
\[
V=\begin{pmatrix}
0.10^2 & 0.2\times 0.10\times 0.20\\
0.2\times 0.10\times 0.20 & 0.20^2
\end{pmatrix}
=
\begin{pmatrix}
0.01&0.004\\0.004&0.04
\end{pmatrix}.
\]
最优线性无偏估计满足
\[
w=\frac{V^{-1}\mathbf 1}{\mathbf 1^T V^{-1}\mathbf 1},
\qquad \mathbf 1=\begin{pmatrix}1\\1\end{pmatrix}.
\]
计算得
\[
w_1=\frac67,\qquad w_2=\frac17.
\]
故
\[
\hat A=\frac67\times 1.20+\frac17\times 1.00
=\frac{8.2}{7}\approx 1.1714.
\]
合并方差为
\[
\sigma_{\hat A}^2=\frac{1}{\mathbf 1^T V^{-1}\mathbf 1}\approx 0.00914,
\]
故
\[
\sigma_{\hat A}\approx 0.0956.
\]
因此
\ansmath{\hat A=1.171\pm 0.096}
\end{solution}
\end{question}

\begin{question}
某模型预言 4 个箱中的期望计数之比为
\[
1:2:1:1.
\]
归一化系数未知。现观测到的四个箱计数分别为
\[
18,\ 35,\ 20,\ 27.
\]
试作 $\chi^2$ 拟合优度检验，并判断是否需要拒绝该模型。可取自由度为 3 时 $5\%$ 显著性水平的临界值
\[
\chi^2_{0.95}(3)=7.815.
\]
\begin{solution}
设期望计数为
\[
E_1=A,\quad E_2=2A,\quad E_3=A,\quad E_4=A.
\]
由总计数
\[
18+35+20+27=100
\]
得
\[
5A=100,
\qquad A=20.
\]
故期望计数为
\[
(20,40,20,20).
\]
于是
\[
\chi^2=\frac{(18-20)^2}{20}+\frac{(35-40)^2}{40}+\frac{(20-20)^2}{20}+\frac{(27-20)^2}{20}.
\]
计算得
\[
\chi^2=\frac4{20}+\frac{25}{40}+0+\frac{49}{20}=3.275.
\]
因为归一化参数是由数据拟合得到的 1 个参数，所以自由度为
\[
4-1=3.
\]
比较可知
\[
3.275<7.815.
\]
故
\anstext{在 $5\%$ 显著性水平下，不拒绝该模型。}
\end{solution}
\end{question}

\begin{question}
做线性刻度实验，模型强制通过原点：
\[
y=bx.
\]
现有四个测量点
\[
(1,1.2),\ (2,2.1),\ (3,2.9),\ (4,4.2),
\]
认为每个 $y$ 的测量标准差都为 $0.1$，彼此独立。用最小二乘法求 $b$、$\sigma_b$，并给出 $\chi^2_{\min}$。
\begin{solution}
由于各点误差相同，最小二乘估计为
\[
\hat b=\frac{\sum x_i y_i}{\sum x_i^2}
=\frac{1\times1.2+2\times2.1+3\times2.9+4\times4.2}{1^2+2^2+3^2+4^2}
=\frac{30.9}{30}=1.03.
\]
故
\ansmath{\hat b=1.03}
参数方差为
\[
\sigma_b^2=\frac{\sigma^2}{\sum x_i^2}=\frac{0.1^2}{30},
\qquad
\sigma_b=\frac{0.1}{\sqrt{30}}\approx 0.0183.
\]
因此
\ansmath{\sigma_b\approx 0.0183}
再算残差平方和：
\[
\chi^2_{\min}=\sum_{i=1}^4\frac{(y_i-\hat b x_i)^2}{0.1^2}.
\]
代入可得
\[
\chi^2_{\min}=7.30.
\]
所以
\ansmath{\chi^2_{\min}=7.30}
\end{solution}
\end{question}

\begin{question}
一次稀有信号搜索中观测到 $n=2$ 个事例，已知本底均值为 $b=0.5$，信号均值记为 $s$。若采用平坦先验
\[
\pi(s)\propto 1,
\qquad s\ge 0,
\]
求 $s$ 的后验众数和后验均值。
\begin{solution}
似然函数为
\[
L(s)\propto e^{-(s+b)}\frac{(s+b)^2}{2!}.
\]
乘以平坦先验后，去掉与 $s$ 无关的常数，得
\[
p(s\mid n=2)\propto e^{-s}(s+0.5)^2,
\qquad s\ge 0.
\]
后验众数由
\[
\frac{d}{ds}\ln p(s\mid n)= -1+\frac{2}{s+0.5}=0
\]
给出，因此
\[
s_{\rm mode}=1.5.
\]
故
\ansmath{s_{\rm mode}=1.5}
后验均值为
\[
E[s\mid n=2]=\frac{\int_0^{\infty} s e^{-s}(s+0.5)^2\,ds}{\int_0^{\infty} e^{-s}(s+0.5)^2\,ds}.
\]
直接积分得
\[
\int_0^{\infty} e^{-s}(s+0.5)^2\,ds=\frac{13}{4},
\qquad
\int_0^{\infty} s e^{-s}(s+0.5)^2\,ds=\frac{33}{4}.
\]
因此
\[
E[s\mid n=2]=\frac{33/4}{13/4}=\frac{33}{13}\approx 2.538.
\]
故
\ansmath{E[s\mid n=2]=\frac{33}{13}\approx 2.538}
\end{solution}
\end{question}

